\documentclass[12pt]{article}
\usepackage[utf8]{inputenc}
\usepackage{amsmath}
\usepackage{amsfonts}
\usepackage{geometry}
\usepackage{enumitem}
\usepackage{parskip}
\usepackage{hyperref}
\usepackage{graphicx}
\graphicspath{ {./images/} }

\geometry{a4paper, margin=1in}

\title{KCPC Workshop 2 \\ Problems: Techniques for Investigation}
\author{King's Competitive Programming Club}
\date{}

\begin{document}

\maketitle

\noindent
Welcome to Workshop 2 of the KCPC 2nd Semester series! This workshop focuses on general mathematical and algorithmic heuristics for investigation: thinking in 3D, wishful thinking, small-case experimentation, induction, visual arguments, recasting problems through invariants and coloring, change of viewpoint, symmetry, and the extreme principle.

\vspace{1em}
\hrule
\vspace{1em}

\section{Psychological Strategies}

\subsection{Pat's Dilemma}
\textbf{Prerequisites:} 3D spatial reasoning, Pythagorean theorem in three dimensions.

Pat wants to take a 1.5-metre-long antique sword onto a train, but the conductor will not allow it as carry-on luggage. The train's baggage master is willing to transport it as checked baggage, but strictly enforces the luggage rule: \textit{no rectangular box or parcel is accepted if any of its individual linear dimensions (length, width, or height) exceeds 1 metre}. The sword is rigid and cannot be bent, disassembled, or warped.

\textbf{Question:} How can Pat package the sword in a rectangular container so that it satisfies the baggage regulations and is accepted for transport?

\subsection{Thinking in 3D}
\textbf{Prerequisites:} Polyhedral geometry, dihedral angles, 3D visualization.

We have two solid polyhedra whose edges all have length 1:
\begin{enumerate}[label=(\roman*)]
    \item A pyramid with a square base of side length 1, whose four triangular faces are equilateral triangles of side length 1.
    \item A regular tetrahedron, which has four equilateral triangular faces of side length 1.
\end{enumerate}
Suppose we glue the regular tetrahedron to the square pyramid along one of the pyramid's triangular faces so that the two triangular faces coincide perfectly.

\textbf{Question:} How many polygonal faces does the resulting combined polyhedron have?

\section{Wishful Thinking and Make It Easier}

\subsection{Long Sequence}
\textbf{Prerequisites:} Pattern recognition, recurrence relations, strong mathematical induction.

The sequence $a_0, a_1, a_2, \ldots$ of real numbers satisfies the relation
\begin{equation}
    a_{m+n} + a_{m-n} = \frac{1}{2}(a_{2m} + a_{2n}) \notag
\end{equation}
for all non-negative integers $m$ and $n$ with $m \ge n$. Given that $a_1 = 1$, find the exact value of $a_{6767}$, and prove your answer rigorously.

\subsection{An \textit{Algorithm?}}
\textbf{Prerequisites:} Binary representations, recursive definitions, bitwise operations.

Let $\mathbb{N} = \{1, 2, 3, 4, \dots\}$ denote the set of positive integers. Consider a function $f: \mathbb{N} \to \mathbb{N}$ that satisfies:
\begin{itemize}
    \item $f(1) = 1$,
    \item $f(2n) = f(n)$ for all $n \in \mathbb{N}$,
    \item $f(2n+1) = f(2n) + 1$ for all $n \in \mathbb{N}$.
\end{itemize}

\textbf{Question:} Find a simple, elegant characterisation of $f(n)$ for any $n \in \mathbb{N}$. Your algorithm should be at most a single sentence long.

\section{Induction}
\textbf{Prerequisites:} Principle of mathematical induction, algebraic inequalities.

\textbf{Problem:} Prove that if $n$ is an integer greater than 3 ($n \ge 4$), then
\begin{equation}
    n! > 2^n. \notag
\end{equation}

\section{Draw a Picture}
\textbf{Prerequisites:} Visual proofs, geometric summation, triangular numbers.

You can prove by algebraic induction that the sum of the first $n$ perfect cubes is equal to the square of the $n$-th triangular number:
\begin{equation}
    \sum_{i=1}^n i^3 = \left( \sum_{i=1}^n i \right)^2 = \frac{n^2(n+1)^2}{4}. \notag
\end{equation}
\textbf{Question:} Can you establish this identity using a geometric picture rather than algebraic induction? Provide a clear diagrammatic argument.

\section{Recast}
\textbf{Prerequisites:} Invariance, parity, bipartite graphs, chessboard colorings.

Two diagonally opposite corner squares are removed from a standard $8 \times 8$ chessboard, leaving a 62-square grid.

\textbf{Question:} Is it possible to tile this mutilated board completely using 31 dominoes of size $2 \times 1$ without overlaps or overhangs? Prove your assertion.

\section{Change Your Point of View}
\textbf{Prerequisites:} Relative motion, angular velocities, frames of reference.

\textbf{Question:} What is the first time after 12 o'clock at which the hour hand and the minute hand of an analog clock overlap exactly? Solve this in your head by choosing a natural, rotating point of view rather than solving long equations.

\section{Symmetry}

\subsection{Grandma's Cabin}
\textbf{Prerequisites:} Geometric reflection principle (Heron's shortest path problem), Euclidean distance.

Your cabin is located 2 miles due north of a straight east-west flowing stream. Your grandmother's cabin is located 12 miles west and 1 mile north of your cabin (so 3 miles north of the stream). Every morning, you walk from your cabin to Grandma's cabin, but first stop at the stream to fill a bucket with fresh water.

\textbf{Question:} What is the minimum total walking distance from your cabin to the stream and then to Grandma's cabin?

\subsection{Geometric Symmetry}
\textbf{Prerequisites:} Rotational symmetry, area scaling, Euclidean geometry.

A smaller square is inscribed in a circle, and that circle is itself inscribed inside a larger square.

\begin{center}
\includegraphics[width=0.28\linewidth]{sym.png}
\end{center}

\textbf{Question:} Find the exact ratio of the area of the smaller square to the area of the larger square using symmetry alone.

\subsection{Symmetry in Probability}
\textbf{Prerequisites:} Continuous/discrete uniform distributions, linearity of expectation, indicator variables.

Imagine dropping 3 pins independently and uniformly at random onto the unit interval $[0, 1]$. They divide the interval into 4 sub-intervals. By symmetry, the expected length of each piece is $1/4$. This can be seen by dropping 4 pins uniformly onto a circle of unit circumference, cutting the circle at the 4th pin, and unrolling it.

Apply similar symmetry arguments to answer the following:
\begin{enumerate}[label=(\alph*)]
    \item An ordinary standard deck of 52 playing cards containing 4 aces is thoroughly shuffled. Cards are dealt one by one from the top until the first ace appears. What is the expected (average) number of cards dealt, including the first ace?
    \item (\textit{Putnam 1992}) Four points are chosen independently and uniformly at random on the surface of a three-dimensional sphere. What is the probability that the centre of the sphere is contained strictly inside the tetrahedron formed by the four points?
\end{enumerate}

\section{The Extreme Principle}
\textbf{Prerequisites:} Proof by contradiction, well-ordering principle, extremal configurations.

Let $B$ and $W$ be finite non-empty sets of black points and white points, respectively, in the Euclidean plane, with the property that every line segment connecting two points of the same colour contains at least one point of the other colour.

\textbf{Problem:} Prove that all points in $B \cup W$ must lie on a single straight line.

\section{Advanced Homework (Competitive Programming)}
\noindent
\textit{Note: These problems are advanced competitive programming challenges selected to reinforce the concepts of invariants, recursive reduction, and parity under contest time and memory constraints.}

\subsection{\href{https://atcoder.jp/contests/arc122/tasks/arc122_c}{ARC 122, Problem C} --- Calculator}
\textbf{Prerequisites:} Fibonacci numbers, Zeckendorf representation / Euclidean-style addition sequences.

Beed has two integers, $x$ and $y$. Initially, $x = 0$ and $y = 0$.

Beed can perform the following four operations any number of times in any order:
\begin{itemize}
    \item Operation 1: $x \leftarrow x + 1$
    \item Operation 2: $y \leftarrow y + 1$
    \item Operation 3: $x \leftarrow x + y$
    \item Operation 4: $y \leftarrow x + y$
\end{itemize}

Given a positive integer $N$ ($1 \le N \le 10^{18}$), output a sequence of at most 130 operations such that the final value of $x$ equals $N$. (The final value of $y$ does not matter.)

\subsection{\href{https://atcoder.jp/contests/arc131/tasks/arc131_c}{ARC 131, Problem C} --- Zero XOR}
\textbf{Prerequisites:} Bitwise XOR properties, game state reduction, parity invariants.

There are $N$ cookies on a desk ($1 \le N \le 400\,000$). Each cookie has a positive integer written on its surface: $A_1, A_2, \dots, A_N$, and all these integers are pairwise distinct ($A_i \neq A_j$ for $i \neq j$).

Two players play a game with these cookies, taking turns alternately. On a turn, a player chooses one cookie from the desk and eats it. If the bitwise XOR sum of the integers written on all remaining cookies on the desk becomes $0$, that player wins immediately, and the game ends.

You go first, and Matthew goes second. Determine whether you have a winning strategy assuming both players play optimally.

\end{document}
